\section{Постановка задач}

Дано обыкновенное дифференциальное уравнение первого порядка с начальным условием вида:
\[
  y'=f(x,y),\quad y(a)=y_0
\]

Необходимо найти приближённое решение задачи Коши на интервале $[a, b]$ с заданным параметром точности $\varepsilon$ --- максимальной допустимой разницей между значениями $y(b)$, которые получены при последовательном уточнении шага интегрирования.

В \textit{варианте №1} вычисление осуществляется \textit{методом Эйлера}. При реализации алгоритма необходимо учитывать следующие условия:
\begin{itemize}
  \item правая часть уравнения $f(x, y)$ выбирается по номеру из набора подготовленных дифференциальных уравнений;
  \item шаг интегрирования $h$ должен определяться и автоматически пересчитываться до тех пор, пока не будет выполнено условие точности $|y_{h/2}(b) - y_h(b)| \leq \varepsilon$;
  \item на графиках должны быть отображены начальное условие $(a, y_0)$, численное решение, полученное методом Эйлера, и аналитическое решение \textit{(при его наличии)} для визуального сравнения и оценки накопления погрешности;
  \item в ходе анализа необходимо исследовать устойчивость метода, а также рассмотреть сценарии практического применения метода.
\end{itemize}

То есть в ходе лабораторной работы будут выполнены следующие задачи:

\begin{enumerate}
  \item Разработка итерационного алгоритма решения задачи Коши методом Эйлера с автоматическим подбором шага $h$ и контролем точности по заданному параметру $\varepsilon$.
  \item Реализация механизма выбора функции $f(x, y)$ по индексу, построения графиков численного и аналитического решений, а также визуализации начальных условий.
  \item Организация ввода данных в формате $a$ $b$ $y_0$ $f_{\text{n}}$ $\varepsilon$, формирование вывода в виде приближённого значения $y(b)$ и проведение анализа устойчивости и погрешности полученного решения.
\end{enumerate}

Теоретический раздел подготовлен с опорой на внешний источник\cite{burden2011numerical}.